\section{Health Endpoints}

Health endpoints are public utility endpoints used to check whether the backend application and database connection are available.

These endpoints do not require authentication.

\subsection{\texttt{GET /health}}

\textbf{Access:} Public

\textbf{Description:} Checks whether the backend application is running.

\textbf{Request body:} None.

\textbf{Example request:}

\begin{minted}[frame=single]{http}
GET /health
\end{minted}

\textbf{Example success response:}

\begin{minted}[frame=single]{json}
{
    "status": "ok",
    "timestamp": "2026-05-30T10:00:00.000Z"
}
\end{minted}

\textbf{Response fields:}

\begin{itemize}
    \item \texttt{status} --- application status. Returns \texttt{ok} when the backend is running.
    \item \texttt{timestamp} --- ISO timestamp generated when the health check response is created.
\end{itemize}

\subsection{\texttt{GET /health/db}}

\textbf{Access:} Public

\textbf{Description:} Checks whether the backend can reach the configured MongoDB database.

\textbf{Request body:} None.

\textbf{Example request:}

\begin{minted}[frame=single]{http}
GET /health/db
\end{minted}

\textbf{Example success response when connected:}

\begin{minted}[frame=single]{json}
{
    "database": "connected",
    "timestamp": "2026-05-30T10:00:00.000Z"
}
\end{minted}

\pagebreak
\textbf{Example response when disconnected:}

\begin{minted}[frame=single]{json}
{
    "database": "disconnected",
    "timestamp": "2026-05-30T10:00:00.000Z"
}
\end{minted}

\textbf{Response fields:}

\begin{itemize}
    \item \texttt{database} --- database connection status. Possible values: \texttt{connected}, \texttt{disconnected}.
    \item \texttt{timestamp} --- ISO timestamp generated when the database health check response is created.
\end{itemize}

\textbf{Important rules:}

\begin{itemize}
    \item This endpoint checks the database connection using a database ping.
    \item The endpoint is public so it can be used by simple monitoring tools or manual development checks.
\end{itemize}